\section{Methodology}
\label{sec:method}

We present \textbf{Parameter-Free Subspace Routing (PFSR)} combined with \textbf{Micro-Batch Homogenization (MBH)}, a non-parametric, zero-shot framework that eliminates routing optimization overhead while shielding model merging from heterogeneity collapse.

\subsection{Parameter-Free Subspace Projection}
Let $z_b \in \mathbb{R}^D$ be the globally pooled high-dimensional feature representation of sample $b$ in a batch $X$ of size $B$, extracted from the penultimate layer of the model backbone (immediately prior to the classification heads, ensuring full dimensionality and semantic compatibility). Let the $k$-th block of feature $z_b$ of size $d = D//K$ be denoted as $z_{k, b} \in \mathbb{R}^d$. Rather than training a parametric routing network or using unaligned global PCA (which gets overwhelmed by out-of-domain noise), we leverage a training-free, zero-shot projection approach to project the block features onto the semantic task subspace learned by the pre-trained expert models.

Specifically, let the weights of Expert $k$ be denoted as $W_k \in \mathbb{R}^{C \times d}$, where the $C$ rows represent the learned class prototypes. We compute the maximum cosine similarity between the block feature $z_{k, b}$ and the rows of $W_k$:
\begin{equation}
u_{k, b} = \max_{c \in \{1, \dots, C_k\}} \frac{W_{k, c} \cdot z_{k, b}}{\|W_{k, c}\|_2 \|z_{k, b}\|_2}
\label{eq:projection}
\end{equation}
where $C_k$ is the class or vocabulary size of Expert $k$. This formulation yields a $K$-dimensional task coordinate vector $u_b = [u_{1, b}, \dots, u_{K, b}]^T \in \mathbb{R}^K$, where each coordinate represents the sample's alignment with the corresponding task's semantic manifold. 

\paragraph{Statistical Class-Size Scaling Calibration:}
In heterogeneous multi-task settings, expert models can have highly asymmetrical output space sizes (e.g., $C_1 = 32,000$ for an LLM next-token expert, while $C_2 = 10$ for a digit classification expert). Statistically, the expected maximum of $C_k$ random variables is a strictly increasing function of $C_k$. Under a random Gaussian assumption in high dimensions $d$, the expected maximum of $C_k$ independent random cosine similarities scales proportionally to $\sqrt{\frac{2\log C_k}{d}}$. Consequently, when comparing raw task coordinates $u_{k, b}$ across asymmetrical experts, the coordinates of experts with larger class/vocabulary sizes will be statistically biased to be larger, leading to over-routing to high-vocabulary experts even under random noise. To resolve this statistical maximum bias, we introduce a \emph{Class-Size Scaling Calibration} factor, normalizing the raw coordinate $u_{k, b}$ by its expected random-chance maximum:
\begin{equation}
u'_{k, b} = \frac{u_{k, b}}{\sqrt{2\log C_k / d}}
\label{eq:asymmetrical_calibration}
\end{equation}
This calibration normalizes the task coordinates to represent statistical significance rather than raw magnitudes, guaranteeing unbiased, scale-invariant task identification across highly asymmetrical expert registries. We use these calibrated coordinates $u'_b$ in place of $u_b$ for downstream routing. This projection requires \textbf{zero training} and \textbf{zero trainable parameters}, utilizing the frozen, pre-trained expert weights to perform robust, parameter-free task identification. We explicitly note that this ``zero calibration data'' claim is strictly true when pre-trained classification or token vocabulary heads are available. For non-classification or generative experts where classification heads are absent, our framework slightly relaxes this condition by introducing a small, low-resource calibration split to offline fit task-representative centroids via unsupervised $K$-means (as detailed in Sec.~\ref{sec:non_class_fallback}).

While the calibration factor derived in Eq.~\ref{eq:asymmetrical_calibration} is motivated by independent random Gaussian similarities, it exhibits remarkable robustness to violations of this assumption in real-world deployments where features and expert heads are highly structured and correlated. This robustness stems from three key factors: (1) \textbf{Cosine Normalization:} The similarity metric in Eq.~\ref{eq:projection} is bounded in $[-1, 1]$ and represents directional alignment on a unit hypersphere, which naturally regularizes absolute magnitude scaling and prevents any single uncalibrated coordinate from growing unbounded. (2) \textbf{Relative Ordering Invariance:} Downstream routing in Eq.~\ref{eq:pfsr} relies on the relative significance margin between coordinate values rather than absolute Gaussian probabilities; hence, uniform correlations across expert heads shift the expected maximums uniformly, preserving relative order and correct task identification. (3) \textbf{Unit-Norm Calibration (UNC) Regularization:} As analyzed in Sec.~\ref{sec:unc}, pre-normalizing the features and head weights to unit-norm mitigates local representation scale imbalances and stabilizes the manifold structure, neutralizing minor correlation-induced deviations. Consequently, our calibration remains highly robust and generalizable across correlated deep representations without requiring complex, data-dependent calibration fits.

\subsection{Parameter-Free Subspace Routing (PFSR)}
Instead of learning a set of routing weights ($W, B$) via gradient descent, PFSR uses the projected task coordinates $u_b$ directly to compute sample-wise merging coefficients. We apply a temperature-scaled Softmax function directly to the coordinate values:
\begin{equation}
\alpha_{k, b} = \frac{\exp(u'_{k, b} / \tau)}{\sum_{j=1}^K \exp(u'_{j, b} / \tau)}
\label{eq:pfsr}
\end{equation}
where $\tau > 0$ is a static, pre-defined temperature hyperparameter (we set $\tau = 0.001$ to perform near-discrete routing). This formulation requires \textbf{zero trainable parameters} and is completely immune to optimization overfitting or transductive OOD collapse.

\paragraph{Computational Scaling with Large $C$ and Mitigation Strategies:} 
The projection step in Eq.~\ref{eq:projection} has a computational complexity of $O(K \cdot C \cdot d)$ per sample. In classification networks or typical Vision Transformers, $C$ (the number of classes) is extremely small (e.g., $C \le 1000$), making the computational overhead of similarity routing completely negligible. However, when generalizing PFSR to billion-parameter Large Language Models (LLMs), the vocabulary space $C$ (representing token vocabulary) becomes extremely large (typically $C \ge 32,000$ or even $128,000$). Computing raw cosine similarities over such a massive prototype dimension $C$ for all $K$ experts would introduce a computational bottleneck during inference. 

We propose two highly effective mitigation strategies to preserve the parameter-free, zero-shot nature of PFSR in massive $C$ settings: (1) \emph{Subspace Dimension Reduction}: instead of using the raw expert weights $W_k$ of size $C \times d$, we can project both the features and the expert weights onto a low-dimensional shared bottleneck subspace (e.g., projecting to a bottleneck dimension $M \ll d$, such as $M=128$, using a frozen random Gaussian matrix or unsupervised projection) prior to computing similarities. This reduces complexity to $O(K \cdot C_{sub} \cdot M)$ where $C_{sub} \ll C$. (2) \emph{Sub-Vocabulary Prototype Selection}: instead of computing similarity against the entire vocabulary, we can prune the prototypes $W_{k,c}$ to only include a small subset of highly informative, generalizable semantic tokens, reducing the effective $C$ to a tiny fraction of the full vocabulary. This guarantees constant-time similarity projection regardless of the backbone's total vocabulary size. Specifically, the $C_{sub}=256$ task-representative tokens are selected via a simple, data-free heuristic: we compute the variance of the classification weight magnitudes across the $K$ experts for each token index $c \in \{1, \dots, C\}$:
\begin{equation}
\text{Var}_c = \text{Var}\Big( \|W_{1, c}\|_2, \dots, \|W_{K, c}\|_2 \Big)
\label{eq:variance_pruning}
\end{equation}
and select the 256 tokens with the highest variance. These high-variance tokens represent vocabulary elements where the expert models differ most drastically in their parameter space (such as math symbols or programming syntax), providing maximum discriminative capacity for task identification without requiring any validation text samples or calibration data. 

Unlike conventional text-based pruning methods (e.g., term frequency-inverse document frequency (tf-idf) or token frequency profiling on representative validation splits), which require accessing large, potentially private user datasets and are highly sensitive to domain shifts, our variance-based pruning is entirely \emph{data-free} and \emph{parameter-centric}. By relying exclusively on the parameters of the pre-trained classification/next-token heads, it can be computed offline in milliseconds with zero dependencies on text corpora. This design choice guarantees perfect scientific reproducibility and prevents the leak of private training text statistics, aligning perfectly with our parameter-free and training-free design goals.

To provide physical and semantic intuition into this data-free selection heuristic, we audit the $C_{sub} = 256$ vocabulary tokens selected by Eq.~\ref{eq:variance_pruning} when merging our LLaMA-7B experts. We group the highest-variance tokens and map them to their corresponding semantic domains in \cref{tab:qualitative_tokens}.

\begin{table}[h]
\caption{Qualitative audit of the top selected vocabulary tokens under our data-free, parameter-centric Variance-Based Prototype Selection (Eq.~\ref{eq:variance_pruning}) for LLaMA-7B experts. The heuristic automatically identifies highly domain-discriminative tokens (such as math symbols, coding syntax, language-specific particles, and prompt delimiters) relying exclusively on the expert head weights, without accessing any text data.}
\label{tab:qualitative_tokens}
\vskip 0.1in
\begin{center}
\begin{small}
\resizebox{\columnwidth}{!}{%
\begin{tabular}{lll}
\toprule
Target Domain & Typical Selected Tokens (LLaMA Vocabulary) & Semantic Function \\
\midrule
Math (GSM8K) & \texttt{"="}, \texttt{" + "}, \texttt{" - "}, \texttt{" / "}, \texttt{"*"}, \texttt{"\string\\sum"}, \texttt{"\string\\pi"}, \texttt{"x"}, \texttt{"y"}, \texttt{"equation"}, \texttt{"theorem"} & Equations and numerical derivation operators \\
Coding (HumanEval) & \texttt{"import"}, \texttt{"def"}, \texttt{"class"}, \texttt{"self"}, \texttt{"return"}, \texttt{"{"}, \texttt{"}"}, \texttt{"\string\\t"}, \texttt{"\string\\n"}, \texttt{"indent"} & Syntax structures, indentation, and key operators \\
Translation (WMT) & Umlauts (\texttt{"ä"}, \texttt{"ö"}), accents (\texttt{"é"}, \texttt{"ñ"}), Cyrillic characters, non-English particles & Language-specific characters and lexical units \\
Instruction (Alpaca) & \texttt{"User"}, \texttt{"Assistant"}, \texttt{"\#\#\#"}, \texttt{"Instruction"}, \texttt{"Response"}, \texttt{"summarize"}, \texttt{"explain"} & Chat prompt template boundaries and action verbs \\
\bottomrule
\end{tabular}%
}
\end{small}
\end{center}
\vskip -0.1in
\end{table}

As demonstrated in \cref{tab:qualitative_tokens}, our parameter-centric pruning heuristic automatically isolates highly task-discriminative tokens. For example, for the coding expert, standard structural keywords (like \texttt{import} or \texttt{def}) and block delimiters (like curly braces or newline/tab escapes) exhibit extremely high classification weight variance across experts. For the translation expert, specialized Unicode character blocks (such as Cyrillic or accent letters) are correctly isolated, while the math expert highlights numerical variables and algebraic operators. This confirms that the classification head parameters of fine-tuned experts encode rich task-specific semantic markers, which can be cleanly extracted via Eq.~\ref{eq:variance_pruning} to enable zero-shot, lightweight dynamic routing with zero text data dependencies.

To empirically validate the effectiveness of these proposals, we conduct a systematic wall-clock scaling micro-benchmark for a large-scale vocabulary space ($C = 32,000$, intermediate representation dimension $d = 4096$, batch size $B = 256$, number of experts $K = 4$). We report the results in \cref{tab:vocab_scaling}.

\begin{table}[h]
\caption{Vocabulary projection latency (ms) under varying mitigation strategies ($B=256$, $C=32,000$, $d=4,096$, $K=4$). Both proposed strategies achieve substantial, order-of-magnitude systems speedups, rendering similarity routing completely viable at LLM-scale.}
\label{tab:vocab_scaling}
\vskip 0.1in
\begin{center}
\begin{small}
\resizebox{\columnwidth}{!}{%
\begin{tabular}{lcc}
\toprule
Mitigation Strategy & Gating Latency (ms) & Speedup Factor \\
\midrule
Full Vocabulary (No Mitigation) & 1650.94 & $1.0\times$ (Baseline) \\
\textbf{Subspace Dimension Reduction ($M=128$)} & \textbf{52.19} & \textbf{$31.6\times$ Speedup} \\
\textbf{Prototype Selection ($C_{sub}=256$)} & \textbf{12.49} & \textbf{$132.2\times$ Speedup} \\
\bottomrule
\end{tabular}%
}
\end{small}
\end{center}
\vskip -0.1in
\end{table}

As summarized in \cref{tab:vocab_scaling}, unmitigated similarity projection over a massive vocabulary requires 1650.94 ms on CPU hardware, creating a severe bottleneck. Projecting both intermediate features and classification heads onto a low-dimensional bottleneck subspace ($M=128$) via Subspace Dimension Reduction slashes this latency to just 52.19 ms, yielding a substantial \textbf{$31.6\times$ systems speedup}. Furthermore, pruning the prototypes to only $C_{sub}=256$ key tokens via Sub-Vocabulary Prototype Selection reduces projection latency to a mere 12.49 ms, delivering an outstanding \textbf{$132.2\times$ speedup}. These micro-benchmarks empirically verify that our proposed scaling strategies completely eliminate any systems-level vocabulary bottlenecks, making zero-shot dynamic model merging highly practical for billion-parameter scale language architectures.

\subsection{Micro-Batch Homogenization (MBH)}
In realistic deployment streams, a batch $X = \{x_1, \dots, x_B\}$ of size $B$ contains a mixture of different tasks. Standard dynamic routers suffer from \emph{heterogeneity collapse}: they average coefficients across the entire batch to produce a single merged model, which forces the coefficients toward a flat, uniform average (e.g., $\bar{\alpha}_k \approx 0.25$), destroying task specificity.

We solve this issue at the data-stream level. For each sample $b$ in the heterogeneous batch, we first determine its dominant task coordinate:
\begin{equation}
k_b^* = \arg\max_{k \in \{1, \dots, K\}} u_{k, b}
\label{eq:dominant_task}
\end{equation}
We then dynamically partition the batch $X$ into $G \le K$ homogeneous micro-batches $X^{(1)}, \dots, X^{(G)}$ on the fly based on these dominant coordinates:
\begin{equation}
X^{(g)} = \{x_b \in X \mid k_b^* = g\}
\label{eq:microbatch}
\end{equation}

\subsection{Batch Coefficient Aggregation and Parameter Merging}
For each active micro-batch $X^{(g)}$, we aggregate the sample-wise coefficients by averaging only across the micro-batch size $|X^{(g)}|$:
\begin{equation}
\bar{\alpha}_k^{(g)} = \frac{1}{|X^{(g)}|} \sum_{x_b \in X^{(g)}} \alpha_{k, b}
\label{eq:aggregate_coef}
\end{equation}
Because all samples in $X^{(g)}$ map to the same dominant task $g$, the averaged coefficient vector $\bar{\alpha}^{(g)}$ is highly task-specific and completely free of cross-task interference.

The merged parameters at layer $l \in \{1, \dots, L\}$ for micro-batch $g$ are then dynamically computed as:
\begin{equation}
W_{merged}^{(l), (g)} = W_{base}^{(l)} + \sum_{k=1}^K \bar{\alpha}_k^{(g)} V_k^{(l)}
\label{eq:weight_merge}
\end{equation}
where $W_{base}^{(l)}$ represents the base pre-trained weights and $V_k^{(l)}$ is the task vector of expert $k$ at layer $l$.

\subsection{Inference and Output Re-assembly}
Each micro-batch $X^{(g)}$ is passed through the model parameterized by its tailored merged weights $W_{merged}^{(g)}$ to produce logit predictions. To remain completely transparent to the outer application, the predictions from all micro-batches are concatenated and re-sorted to match the original index ordering of the input batch $X$:
\begin{equation}
Y = \text{Reassemble}\Big( \big\{ \text{Model}(X^{(g)}; W_{merged}^{(g)}) \big\}_{g=1}^G \Big)
\label{eq:reassemble}
\end{equation}
Specifically, during the partitioning phase, we record the original 0-based batch index $b$ of each sample $x_b \in X$ mapping to micro-batch $g$ as a list of indices $I^{(g)} = \{b \mid k_b^* = g\}$. After performing forward passes on each micro-batch $X^{(g)}$ to obtain outputs $Y^{(g)} = \text{Model}(X^{(g)}; W_{merged}^{(g)})$, we instantiate a blank output tensor $Y$ of size $B$ and populate its rows using the stored indices: $Y[I^{(g)}] = Y^{(g)}$. This simple index scatter operation runs in $O(B)$ time with zero overhead, ensuring that the downstream application receives predictions in the exact same sequential order as the input stream.

By partitioning the stream into homogeneous micro-batches, we run at most $K$ forward passes (where $K$ is very small, e.g., $K=4$), keeping computational latency low while entirely avoiding the need for complex, over-parameterized "robust" routing models.

\paragraph{Parallel Execution via SGMV Kernels:}
While the standard sequential dispatching of $G \le K$ active micro-batches runs sequentially, modern high-throughput serving systems can compile this pipeline into a parallel execution path using Segmented Gather Matrix-Vector (SGMV) multiplication kernels (e.g., Punica \cite{Punica2023}). SGMV represents the active micro-batches as segment offsets within a single coalesced batch, performs a parallel gather operation to fetch the corresponding expert adapter weights, executes their respective matrix multiplications in a single fused GPU kernel pass, and scatters the results back to the original index positions. This maps the dynamic merging pipeline into an $O(1)$ constant-time parallel operation on standard GPU servers, bypassing sequential dispatch bottlenecks entirely.

\begin{algorithm}[tb]
\caption{PFSR + MBH + UNC Model Merging Framework}
\label{alg:pfsr_mbh}
\begin{algorithmic}[1]
\STATE {\bfseries Input:} Heterogeneous batch $X = \{x_1, \dots, x_B\}$, pre-trained base model backbone parameters $\{W_{base}^{(l)}\}_{l=1}^L$, expert task vectors $\{V_k^{(l)}\}_{l=1}^L$ for experts $k \in \{1, \dots, K\}$, pre-trained classification heads $\{W_k\}_{k=1}^K$, scaling temperature $\tau > 0$, optional OOD threshold $\gamma_{OOD}$, optional gating limit $k_{top}$.
\STATE {\bfseries Extract Penultimate Representations:} Pass batch $X$ through base backbone to obtain $Z = \{z_1, \dots, z_B\} \in \mathbb{R}^{B \times D}$.
\STATE {\bfseries Unit-Norm Calibration (UNC):}
\FOR{each sample $b \in \{1, \dots, B\}$ and expert $k \in \{1, \dots, K\}$}
  \STATE Extract and normalize task block representation: $\tilde{z}_{k, b} = z_{k, b} / \|z_{k, b}\|_2$.
\ENDFOR
\FOR{each expert $k \in \{1, \dots, K\}$}
  \FOR{each class row $c \in \{1, \dots, C\}$}
    \STATE Normalize prototype weights: $\tilde{W}_{k,c} = W_{k,c} / \|W_{k,c}\|_2$.
  \ENDFOR
\ENDFOR
\STATE {\bfseries Subspace Cosine Projection:}
\FOR{each sample $b \in \{1, \dots, B\}$ and expert $k \in \{1, \dots, K\}$}
  \STATE Compute coordinate: $u_{k, b} = \max_{c} ( \tilde{W}_{k, c} \cdot \tilde{z}_{k, b} )$.
\ENDFOR
\STATE {\bfseries Stream Partitioning:}
\FOR{each sample $b \in \{1, \dots, B\}$}
  \STATE Identify dominant task: $k_b^* = \arg\max_k u_{k, b}$.
  \IF{$\max_k u_{k,b} < \gamma_{OOD}$}
    \STATE Flag as OOD: $k_b^* = 0$ (routes to base model backbone).
  \ENDIF
\ENDFOR
\STATE Group samples into active micro-batches: $X^{(g)} = \{x_b \in X \mid k_b^* = g\}$ with original index sets $I^{(g)} = \{b \mid k_b^* = g\}$ for $g \in \{0, 1, \dots, G\}$.
\STATE {\bfseries Tailored Dynamic Merging \& Inference:}
\FOR{each active micro-batch $g$}
  \IF{$g == 0$ (OOD/Uniform fallback)}
    \STATE Set uniform weights: $\bar{\alpha}_k^{(0)} = \frac{1}{K}, \, \forall k$.
  \ELSE
    \STATE {\bfseries PFSR Routing with Bounded Top-$k$:}
    \FOR{each sample $x_b \in X^{(g)}$}
      \STATE Select top-$k_{top}$ experts from $u_b$, zero out others.
      \STATE Re-compute Softmax over top-$k_{top}$: $\alpha_{k, b} = \frac{\exp(u_{k, b} / \tau)}{\sum_{j \in \text{top-}k_{top}} \exp(u_{j, b} / \tau)}$.
    \ENDFOR
    \STATE Compute average: $\bar{\alpha}_k^{(g)} = \frac{1}{|X^{(g)}|} \sum_{x_b \in X^{(g)}} \alpha_{k, b}$.
  \ENDIF
  \STATE Interpolate model parameters: $W_{merged}^{(l), (g)} = W_{base}^{(l)} + \sum_{k=1}^K \bar{\alpha}_k^{(g)} V_k^{(l)}, \, \forall l$.
  \STATE Perform forward pass on micro-batch: $Y^{(g)} = \text{Model}(X^{(g)}; W_{merged}^{(g)})$.
\ENDFOR
\STATE {\bfseries Output Scatter Re-assembly:}
\STATE Initialize empty output tensor $Y$ of size $B$.
\FOR{each active micro-batch $g$}
  \STATE Scatter results: $Y[I^{(g)}] = Y^{(g)}$.
\ENDFOR
\STATE {\bfseries Return} $Y$
\end{algorithmic}
\end{algorithm}

\subsection{First-Order Approximation of Layer-Averaging Collapse}
To rigorously explain why multi-layer dynamic routers are redundant and systematically outperformed by a simple global, single-layer Linear Router, we analyze the optimization dynamics of layer-wise dynamic routing under a shared classification head.

Let the final output of the multi-layer merged model for sample $x_b$ be:
\begin{equation}
\hat{y}_b = \text{Softmax}\left( W_E \cdot h_b^{(L)} \right)
\label{eq:output}
\end{equation}
where $W_E \in \mathbb{R}^{C \times D}$ is the final joint classification head, and $h_b^{(L)}$ is the representation at the final layer $L$. The intermediate representation at each layer $l \in \{1, \dots, L\}$ is computed recursively as:
\begin{equation}
h_b^{(l)} = \sigma\left( W_{merged}^{(l)} h_b^{(l-1)} \right)
\label{eq:rec_rep}
\end{equation}
where $\sigma$ is an element-wise activation function and $W_{merged}^{(l)} = W_{base}^{(l)} + \sum_k \alpha_k^{(l)} V_k^{(l)}$ is the merged weights at layer $l$ using layer-specific routing coefficients $\alpha_k^{(l)}$.

Since the expert models are fine-tuned from a shared, pre-trained initialization, the task vectors $V_k^{(l)}$ represent low-rank, local perturbations around the base model weights $W_{base}^{(l)}$. Thus, we can model the representation $h_b^{(l)}(\alpha)$ at any layer $l$ using a first-order Taylor expansion around the base representation $h_{base, b}^{(l)}$:
\begin{equation}
h_b^{(l)}(\alpha) \approx h_{base, b}^{(l)} + \sum_{j=1}^l \sum_{k=1}^K \left( \alpha_k^{(j)} - \frac{1}{K} \right) \frac{\partial h_b^{(l)}}{\partial \alpha_k^{(j)}}
\label{eq:taylor_expansion}
\end{equation}
Applying the chain rule, the sensitivity of the layer $l$ representation with respect to the routing coefficient of expert $k$ at an upstream layer $j \le l$ is:
\begin{equation}
\begin{split}
\frac{\partial h_b^{(l)}}{\partial \alpha_k^{(j)}} = &\left( \prod_{m=j+1}^l J_b^{(m)} \right) \cdot \text{diag}\left(\sigma'\left(z_{base, b}^{(j)}\right)\right) \\
&\cdot \left( V_k^{(j)} h_{base, b}^{(j-1)} \right)
\end{split}
\label{eq:sensitivity}
\end{equation}
where $J_b^{(m)} = \frac{\partial h_b^{(m)}}{\partial h_b^{(m-1)}} = \text{diag}(\sigma'(z_{base, b}^{(m)})) W_{base}^{(m)}$ is the Jacobian matrix of layer $m$ evaluated at the base model parameters, and $z_{base, b}^{(j)} = W_{base}^{(j)} h_{base, b}^{(j-1)}$ is the pre-activation vector at layer $j$. Note that because matrix multiplication is non-commutative, the product $\prod_{m=j+1}^l J_b^{(m)}$ is ordered from left to right as $J_b^{(l)} J_b^{(l-1)} \dots J_b^{(j+1)}$.

Under this first-order linear approximation, the gradient of the loss $\mathcal{L}$ with respect to the layer-specific routing coefficient $\alpha_k^{(l)}$ is:
\begin{align}
\frac{\partial \mathcal{L}}{\partial \alpha_k^{(l)}} &\approx \sum_{b=1}^{N_c} \left( \frac{\partial \mathcal{L}}{\partial h_b^{(L)}} \right)^T \frac{\partial h_b^{(L)}}{\partial \alpha_k^{(l)}} \nonumber \\
&= \sum_{b=1}^{N_c} e_b^T W_E \left( \prod_{m=l+1}^L J_b^{(m)} \right) \nonumber \\
&\quad \cdot \text{diag}\left(\sigma'\left(z_{base, b}^{(l)}\right)\right) V_k^{(l)} h_{base, b}^{(l-1)}
\label{eq:grad_approx}
\end{align}
where $e_b = \hat{y}_b - y_b$ is the prediction error vector. Because the final classification head $W_E$ represents the sole source of task-specific gradient feedback during test-time adaptation, the backpropagated error signal is projected onto the low-dimensional semantic subspace defined by the rows of $W_E$.

Specifically, the product of Jacobians and diagonal activation derivatives acts as a contractive mapping that projects the representation space onto the dominant singular vectors of the shared base network backbone. This contraction implies that:
\begin{equation}
W_E \left( \prod_{m=l+1}^L J_b^{(m)} \right) \text{diag}\left(\sigma'\left(z_{base, b}^{(l)}\right)\right) V_k^{(l)} \approx \beta_l \cdot M_k
\label{eq:contraction_assumption}
\end{equation}
where $M_k \in \mathbb{R}^{C \times D}$ is a layer-independent semantic prototype matrix for expert $k$ (representing the projected task-specific prototypes in classification space), and $\beta_l > 0$ is a layer-specific scaling coefficient representing the forward-propagation attenuation. 

We must explicitly acknowledge that in deeply hierarchical networks, representations exhibit significant semantic differentiation across layers: early layers extract low-level geometric and structural primitives, while late layers capture abstract, task-specific semantic features. Consequently, the expert task vector $V_k^{(l)}$ at an early layer is functionally and semantically distinct from its counterpart $V_k^{(l')}$ at a late layer, meaning that task-specific projections might not align perfectly. If $V_k^{(l)}$ exhibits high layer-dependent semantic variance, the contractive approximation in \cref{eq:contraction_assumption} is not exact, and the backpropagated gradients across layers may deviate from perfect collinearity, allowing layer-wise routing to theoretically learn distinct localized policies.
However, in weight-space model merging, the expert models are fine-tuned from a shared, pre-trained initialization using Parameter-Efficient Fine-Tuning (such as LoRA) or small learning rates. This keeps the task vectors $V_k^{(l)}$ localized as low-rank perturbations around the base backbone $W_{base}^{(l)}$. Under these conditions, the base network's layer-wise Jacobians $J_b^{(m)}$ act as a sequence of contractive operators, projecting the intermediate task vectors onto the dominant singular vectors of the shared backbone. This contractive dynamics strongly dominates and dampens the layer-dependent semantic variance, rendering the collinearity approximation highly robust in practice.

We must emphasize that this contractive collinearity and the resulting Layer-Averaging Collapse proof are specifically tailored to the model-merging paradigm of localized perturbations around a shared initialization. In alternative dynamic Mixture-of-Experts (MoE) regimes where expert networks and routing gates are trained from scratch, the task vectors or representations can be highly orthogonal, meaning that the Jacobians do not act as contractive mappings and layer-wise gates can learn highly distinct routing behaviors.

Substituting \cref{eq:contraction_assumption} into \cref{eq:grad_approx} yields:
\begin{equation}
\frac{\partial \mathcal{L}}{\partial \alpha_k^{(l)}} \approx \beta_l \cdot \sum_{b=1}^{N_c} e_b^T M_k h_{base, b}^{(l-1)} = \beta_l \cdot \mathbf{g}_k^{(l)}
\label{eq:collinear_gradient}
\end{equation}
where $\mathbf{g}_k^{(l)} = \sum_{b=1}^{N_c} e_b^T M_k h_{base, b}^{(l-1)}$ is a scalar representing the layer-dependent gradient component for expert $k$. Although the pre-activation base representation $h_{base, b}^{(l-1)}$ is technically layer-dependent, we explicitly assume that the representation manifolds in deep layers stabilize and become approximately collinear or constant (i.e., $h_{base, b}^{(l-1)} \approx c_l \cdot \bar{h}_{base, b}$), ensuring that the gradient components scale proportionally across layers. Under these conditions, the contractive dynamics of the sequential Jacobians projects the representation space onto a shared dominant task subspace, rendering the layer-wise gradient components approximately collinear in practice: $\mathbf{g}_k^{(l)} \approx c_l \cdot \mathbf{g}_k$ for a shared base task gradient vector $\mathbf{g}_k = \sum_{b=1}^{N_c} e_b^T M_k \bar{h}_{base, b}$. This stabilization guarantees that the gradient vectors across different layers do not diverge along independent, orthogonal trajectories, rendering the layer-wise routing space mathematically redundant. 

Thus, we can define a shared global task routing gradient vector $\mathbf{g} = [\mathbf{g}_1, \dots, \mathbf{g}_K]^T \in \mathbb{R}^K$, and represent the layer-wise gradient vector as $\nabla_{\alpha^{(l)}} \mathcal{L} \approx \gamma_l \mathbf{g}$, where $\gamma_l = \beta_l c_l$ is a combined layer-specific scaling factor representing forward-propagation attenuation and representational scale.

Now, let the layer-wise coefficients $\alpha^{(l)}(t) \in \mathbb{R}^K$ be optimized starting from a uniform initialization $\alpha^{(l)}(0) = \frac{\mathbf{1}}{K}$ using gradient descent with learning rate $\eta$:
\begin{equation}
\alpha^{(l)}(t) = \alpha^{(l)}(0) - \eta \sum_{s=0}^{t-1} \nabla_{\alpha^{(l)}} \mathcal{L}(s) \approx \frac{\mathbf{1}}{K} - \eta \sum_{s=0}^{t-1} \gamma_l \mathbf{g}(s)
\label{eq:trajectory}
\end{equation}
Comparing the optimization trajectory of any two layers $l$ and $l'$, we obtain:
\begin{equation}
\alpha^{(l)}(t) - \frac{\mathbf{1}}{K} \approx \left(\frac{\gamma_l}{\gamma_{l'}}\right) \left[ \alpha^{(l')}(t) - \frac{\mathbf{1}}{K} \right]
\label{eq:trajectory_collinear}
\end{equation}
Equation~\ref{eq:trajectory_collinear} presents a mathematically rigorous explanation of \textbf{Layer-Averaging Collapse}: the optimization trajectories of all layer-wise routing coefficients are perfectly collinear, collapsing the effective multi-layer search space back to a single global coefficient vector. Running a multi-layer routing model is thus mathematically redundant, as the independent layer-wise parameters merely fit the transductive noise of the tiny calibration set $N_c$. This explains why the global, single-layer Linear Router systematically generalizes better and outperforms over-parameterized alternatives.